\section{Controlling the layer errors}\label{sec:05-layers}

We now bound the terms in Equation~\eqref{eq:18} that still involve
the nonlinear field $A$, its remainder $R=A-L$, and the layer defects
$E_z^L$. Our objective is to bound $\mathcal E$, defined at the end of
Section~\ref{sec:03-entropy}, through functionals of $L$, retaining a
nonpositive integral against the layer measure $\beta=\beta^L$.
Throughout this section, all the constants are those fixed earlier;
in particular, \(t_\pm=t_c\pm t_r\) and \(\Sigma=I+T\).
For a symmetric matrix field \(E\), write
\(\|E\|_\Sigma^2=\tau_\Sigma E^2\).

\subsection{The Gaussian remainder}

The deterministic weight \(\Sigma\) preserves orthogonality of
Gaussian Hermite degrees.  Write
\[
 R_e(G)=\frac{R(G)+R(-G)}2,\qquad
 R_o(G)=\frac{R(G)-R(-G)}2.
\]
The normalization \(\mathbb E A=0\) and the definition of \(L\) as the
degree-one part of \(A\) show that \(R_e\) has degrees at least \(2\),
whereas \(R_o\) has degrees at least \(3\).

\begin{lemma}[Control of the Gaussian remainder]\label{lem:layers-gaussian}
With the notation of Section~\ref{sec:04-quadratic},
\begin{equation}
 \sigma_\Sigma(F)-D_\Sigma(u)^2
 \le (b+\eta)\tau_\Sigma R^2-\eta\tau_\Sigma R_o^2
       +b_m\|M_{f_q}^L\|_2^2.
 \tag{19}\label{eq:19}
\end{equation}
Moreover,
\[
 \|M_{f_q}^L\|_2^2
 \le 2\mathcal L_I[S_q(z)S_q(x)].
\]
\end{lemma}

\begin{proof}
The definitions of \(h_l,f_l,d_l^*\) depend linearly on \(l\).
For the constant profile \(l=1\), its centered primitive is \(x\)
and its image under \(1+\mathcal N\) is \(2x\).  Hence
\[
 d_1^*=2\mathcal K_G A-L=2\mathcal K_G R,
\]
because \(\mathcal K_G L=L/2\).
Since \(F=c+2kq\), it follows that
\[
 \sigma_\Sigma(F)
 =4c\langle R,\mathcal K_G R\rangle_\Sigma
       +4k\langle R,d_q^*\rangle_\Sigma.
\]
The vector \(u=(\pi,q)\) gives
\(D_\Sigma(u)^2\ge\|d_q^*\|_{1,\Sigma}^2\).
We estimate the resulting quadratic expression separately on odd
and even Hermite degrees.

On the odd degrees, completion of the square in the
\((1+\mathcal N_G)\)-form norm gives
\[
 4k\langle R_o,(d_q^*)_o\rangle_\Sigma
       -\|(d_q^*)_o\|_{1,\Sigma}^2
 \le 4k^2\langle R_o,\mathcal K_G R_o\rangle_\Sigma.
\]
Here the square is centered at \(2k\mathcal K_G R_o\), so the
calculation remains valid in the form domain.  Since
\(c+k^2=b\) and \(\mathcal K_G\le1/4\) on the degrees of \(R_o\),
the odd contribution is at most \(b\tau_\Sigma R_o^2\).

For the even degrees, the evenness of \(q\) implies that its
Gaussian-centered primitive \(h_q\) is odd, and hence \(f_q\) is
odd.  The field \(L\) is odd in \(G\), so \(h_q(L)\) and
\(f_q(L)\) are odd matrix fields.  Thus
\[
 (d_q^*)_e=\mathcal K_G(M_{f_q}^L)_e.
\]
Put \(\mathfrak b=b+\eta\) and
\[
 b_m'=\frac{b_m}{1+t_+}
      =\frac{b-3\eta}{3(3\mathfrak b-4c)}=\frac{134}{309}>0.
\]
On a positive even degree, the eigenvalue \(\rho\) of
\(\mathcal K_G\) lies in \((0,1/3]\).
Writing \(R_j,M_j\) for the components of \(R,M_{f_q}^L\) in that degree,
the quadratic form to estimate is
\[
 4c\rho\|R_j\|_\Sigma^2
       +4k\rho\langle R_j,M_j\rangle_\Sigma-\rho\|M_j\|_\Sigma^2.
\]
The scalar matrix inequality
\[
 \begin{pmatrix}4c\rho&2k\rho\\2k\rho&-\rho\end{pmatrix}
 \le
 \begin{pmatrix}\mathfrak b&0\\0&b_m'\end{pmatrix}
\]
holds for \(0\le\rho\le1/3\).  Indeed, the first diagonal entry
of the difference is at least
\[
 \mathfrak b-\frac{4c}{3}=\frac{103}{750}>0,
\]
and its determinant is
\[
 (\mathfrak b-4c\rho)(b_m'+\rho)-4k^2\rho^2
 =(1-3\rho)\left(\mathfrak b\,b_m'+\frac{4b}{3}\rho\right)\ge0.
\]
The constant degree has zero \(R\)-coefficient and contributes a
nonpositive term.  Summing the even degrees therefore gives at most
\[
 \mathfrak b\tau_\Sigma R_e^2
       +b_m'\|(M_{f_q}^L)_e\|_\Sigma^2
 \le \mathfrak b\tau_\Sigma R_e^2+b_m\|M_{f_q}^L\|_2^2,
\]
where we used \(\Sigma\le(1+t_+)I\).
Together with the odd estimate, this proves Equation~\eqref{eq:19}.

Finally, Equation~\eqref{eq:15} yields
\[
 \|M_{f_q}^L\|_2^2
 =2e(f_q,f_q)-\operatorname{Nt}(f_q,f_q).
\]
The kernel defining \(\operatorname{Nt}\) is nonnegative, and
\(f_q'=S_q>0\); hence its last term is nonnegative.
The identity \(e(f_q,f_q)=\mathcal L_I[S_q(z)S_q(x)]\)
proves the remaining assertion.
\end{proof}

\subsection{A positivity fact for spectral layers}

The weight \(T\) need not commute with \(A\), \(L\), or the
projection derivatives \(P_z\).  The following elementary fact
allows us to retain the positive parts of the layer terms without
a commutation assumption.

\begin{lemma}[Positive blocks and weighted diagonal bounds]
\label{lem:layers-diagonal}
Let \(B_*\) be a real symmetric matrix, let \(0\le P\le I\), and
fix \(z\) outside the spectrum of \(B_*\).
In an orthonormal eigenbasis of \(B_*\), write
\[
 E=P-1_{(-\infty,z]}(B_*),\qquad
 a_i=|x_i-z|,\qquad \epsilon_i=\operatorname{sgn}(x_i-z).
\]
Then
\[
 (E^2)_{ii}\le\epsilon_iE_{ii}.
\]
The matrix \(S\) with entries
\[
 S_{ij}=
 \begin{cases}
  \epsilon_i\sqrt{a_i a_j}\,E_{ij},&\epsilon_i=\epsilon_j,\\
  0,&\epsilon_i\ne\epsilon_j
 \end{cases}
\]
is positive semidefinite, and
\(\operatorname{tr}S=\sum_i a_i\epsilon_iE_{ii}\).
For any nonnegative numbers \(w_i\),
\[
 \sum_{i,j}\frac{a_iw_i+a_jw_j}{2}|E_{ij}|^2
 \le\sum_i a_iw_i\epsilon_iE_{ii}.
\]
\end{lemma}

\begin{proof}
Write \(Q=1_{(-\infty,z]}(B_*)\).
Since \(Q\) is diagonal in the chosen basis and \(P^2\le P\),
\[
 (E^2)_{ii}=(P^2)_{ii}-2Q_{ii}P_{ii}+Q_{ii}
 \le P_{ii}-2Q_{ii}P_{ii}+Q_{ii}
 =\epsilon_iE_{ii}.
\]
On the block \(x_i>z\), the compression of \(E\) is the
compression of \(P\), and is positive semidefinite.
On the block \(x_i<z\), the compression of \(-E\) is the
compression of \(I-P\), and has the same property.
Diagonal congruence by the factors \(\sqrt{a_i}\) proves the
assertion about \(S\).
Finally, symmetry of \(E\) gives
\[
 \sum_{i,j}\frac{a_iw_i+a_jw_j}{2}|E_{ij}|^2
 =\sum_i a_iw_i(E^2)_{ii}
 \le\sum_i a_iw_i\epsilon_iE_{ii}.
\]
\end{proof}

We apply Lemma~\ref{lem:layers-diagonal} pointwise in \(G\), with \(P=P_z\) and
\(B_*=A\) or \(L\), and then integrate over \(z\).
The normalized expected integral of
\(\sum_i a_i\epsilon_i(E_z^{B_*})_{ii}\) is
\(\mathcal B^{B_*}\), and its version with a weight \(w(z,x_i)\)
is \(\int w\,d\beta^{B_*}\).
For each \(G\), the excluded spectral endpoints form a finite
set of layers and thus do not affect these integrals.
All subsequent sums use the normalization \(m^{-1}\mathbb E\);
the moment and layer convergence bounds established earlier
justify their integration.

\begin{lemma}[The weighted constant layer]\label{lem:layers-constant}
The weighted non-nesting term satisfies
\begin{equation}
 N_T=2e_T^A(\iota,\iota)
 \ge t_-N-\sqrt{c_gN/2}\,\|[A,T]\|_2,
 \qquad c_g=2\log2-\frac13.
 \tag{20}\label{eq:20}
\end{equation}
\end{lemma}

\begin{proof}
Equation~\eqref{eq:15}, applied with threshold \(A\), gives
\(N_T=2e_T^A(\iota,\iota)\) and \(N=2\mathcal B^A\).
In a spectral basis of \(A\), the scalar kernel multiplying
\((E_z^A)_{ij}T_{ji}\) in \(2e_T^A(\iota,\iota)\) is
\(x_i+x_j-2z\).
Separate this kernel into a comparison term and a remainder. The
comparison term is \(2\epsilon\sqrt{a_i a_j}\) when both endpoints
have sign \(\epsilon\), and zero when they have opposite signs.
Its pairing is \(2\operatorname{tr}(TS_z)\), where \(S_z\ge0\)
is supplied by Lemma~\ref{lem:layers-diagonal}. Thus the original
pairing equals \(2\operatorname{tr}(TS_z)\) plus the pairing with
the remaining kernel. Since \(T\ge t_-I\), the normalized expected
integral of the comparison term is at least
\(2t_-\mathcal B^A=t_-N\).

Let \(k_z(x_i,x_j)\) be the remaining kernel, and put
\(\omega_z=(a_i+a_j)/2\).
On a common-sign block it is
\(\epsilon(\sqrt{a_i}-\sqrt{a_j})^2\);
between the endpoints it remains \(x_i+x_j-2z\).
We claim that
\[
 \int_{\mathbb R}\frac{k_z(x_i,x_j)^2}{\omega_z}\,dz
 =c_g(x_i-x_j)^2.
\]
If the endpoints coincide, the remaining kernel vanishes;
the quotient is assigned value zero wherever its denominator
vanishes.  Otherwise, translation and scaling reduce the
calculation to endpoints \(0,1\).
The interval between them contributes
\[
 2\int_0^1(1-2z)^2\,dz=\frac23.
\]
The two exterior intervals contribute
\[
 4\int_0^\infty
       \frac{(\sqrt{1+t}-\sqrt t)^4}{1+2t}\,dt=2\log2-1.
\]
For completeness, use the scalar substitution
\[
 u_0=(\sqrt{1+t}-\sqrt t)^2,\qquad
 t=\frac{(1-u_0)^2}{4u_0},\qquad
 |dt|=\frac{1-u_0^2}{4u_0^2}\,du_0.
\]
The last integral becomes
\[
 2\int_0^1\frac{u_0(1-u_0^2)}{1+u_0^2}\,du_0
 =2\log2-1.
\]
This proves the claim.

Cauchy--Schwarz in the variables \(G,z,i,j\), with weight
\(\omega_z\), bounds the absolute value of the remaining pairing by
\[
 \left(\frac1m\mathbb E\int
          \sum_{i,j}\omega_z|(E_z^A)_{ij}|^2\,dz\right)^{1/2}
 \left(\frac1m\mathbb E\sum_{i,j}|T_{ji}|^2
          \int\frac{k_z(x_i,x_j)^2}{\omega_z}\,dz\right)^{1/2}.
\]
Lemma~\ref{lem:layers-diagonal} bounds the first factor by
\(\sqrt{\mathcal B^A}\).
The second factor is \(\sqrt{c_g}\|[A,T]\|_2\), since
the commutator has entries \((x_i-x_j)T_{ij}\).
This proves Equation~\eqref{eq:20}.
\end{proof}

\subsection{Combining the quadratic and constant-layer terms}

We have controlled the Gaussian remainder and the constant part of
the \(T\)-weighted layer.  We next apply those estimates to
Equation~\eqref{eq:18}; the resulting bound will leave just the
variable profile \(H_0\) to estimate.

The identity \(H+v''=2\kappa+H_0\) gives
\[
 -2\kappa\mathcal L_T[1]
 =-2\kappa e_T(\iota,\iota)
 =-\kappa\tau_T R^2-\kappa N_T
\]
by Equation~\eqref{eq:15}.
To track the signs when combining this with
Equation~\eqref{eq:19}, put
\[
 a_0=b+\eta+(b+\eta-\kappa)t_-.
\]
Here \(b+\eta-\kappa=-.536<0\), while \(t_-=-.022<0\);
consequently \(a_0\ge b+\eta\).
Since \(R^2\ge0\) and \(T\ge t_-I\),
\[
 \begin{aligned}
 (b+\eta)\tau_\Sigma R^2-\kappa\tau_T R^2
 &=(b+\eta)\tau R^2+(b+\eta-\kappa)\tau_T R^2\\
 &\le a_0\tau R^2
   =2a_0\mathcal B-a_0N\\
 &\le 2a_0\mathcal B-(b+\eta)N.
 \end{aligned}
\]
The equality uses Equation~\eqref{eq:15}, and the last inequality
uses \(N=2\mathcal B^A\ge0\).

Including the term \(-N/8\) in \(\mathcal E\), the remaining
terms from Equation~\eqref{eq:20} are
\[
 -\Lambda_0N+\kappa\sqrt{c_gN/2}\,\|[A,T]\|_2,\qquad
 \Lambda_0=b+\eta+\kappa t_-+\frac18=.72348>0.
\]
Completing the square in \(\sqrt N\) bounds this expression by
\[
 \frac{\kappa^2c_g}{8\Lambda_0}\|[A,T]\|_2^2
 \le e_0\|[A,T]\|_2^2.
\]
This last comparison has a positive margin: using \(\log2<.694\),
\[
 \frac{\kappa^2c_g}{8\Lambda_0}
 <\frac{665231}{2713050}<.245197,
 \qquad e_0=.256.
\]
The elementary bound on \(\log2\), if needed, follows from
\[
 \log2=2\sum_{j=0}^\infty
              \frac1{(2j+1)3^{2j+1}}
 \le \frac23+\frac2{81}+\frac2{1215}+\frac1{6804}<.694.
\]

The fields \(L\) and \(R\) occupy orthogonal Gaussian degrees.
As \(T\) is deterministic, their commutators with \(T\) remain
orthogonal.  Also, in a spectral basis of \(T\), every difference
of two eigenvalues has absolute value at most \(2t_r\).
It follows that
\[
 \begin{aligned}
 e_0\|[A,T]\|_2^2
 &=e_0\|[L,T]\|_2^2+e_0\|[R,T]\|_2^2\\
 &\le e_0\|[L,T]\|_2^2+e_0(2t_r)^2\tau R^2\\
 &\le e_0\|[L,T]\|_2^2+2e_0(2t_r)^2\mathcal B.
 \end{aligned}
\]
Define the scalar profile
\[
 P_0(z,x)=2a_0+2b_mS_q(z)S_q(x)+2e_0(2t_r)^2.
\]
The bound in Lemma~\ref{lem:layers-gaussian} for
\(\|M_{f_q}^L\|_2^2\), together with the preceding inequalities,
now gives
\begin{equation}
 \begin{aligned}
 \mathcal E\le{}&
 \Delta_{\rm p}-J_\Sigma(u)+e_0\|[L,T]\|_2^2
 -\eta\tau_\Sigma R_o^2-\langle R,v'(L)\rangle_T
 +\tau(T-\log\Sigma)\\
 &-\mathcal L_I[H-P_0]-\mathcal L_T[H_0].
 \end{aligned}
 \tag{21}\label{eq:21}
\end{equation}

\subsection{The variable weighted layer}

For a profile \(\Phi(z,x)\), let
\[
 \overline{\Phi}(z,x)=
 \begin{cases}
  \displaystyle\frac1{x-z}\int_z^x\Phi(z,y)\,dy,&x\ne z,\\[2mm]
  \Phi(z,z),&x=z.
 \end{cases}
\]
Thus the bar is the uniform average over the interval with
endpoints \(z,x\), regardless of their order.
Define
\[
 \begin{aligned}
 g_z(x)&=\overline{H_0}(z,x),\\
 W(z,x)&=\overline{H-P_0}(z,x)+t_cg_z(x)
       -\frac{t_r}{2}\left(h_s+\frac{g_z(x)^2}{h_s}\right),\\
 W_0(z,x)&=\frac{\overline{H_0^2}(z,x)}{\alpha}.
 \end{aligned}
\]
Averaging the first inequality in Equation~\eqref{eq:3} gives
\[
 \overline{H-P_0}+t_cg_z
 -\frac{t_r}{2}\left(h_s+\frac{\overline{H_0^2}}{h_s}\right)
 >\frac{\overline{H_0^2}}{\alpha}.
\]
Since \(g_z^2\le\overline{H_0^2}\), we conclude that
\(W>W_0\ge0\), also for coincident endpoints.

\begin{lemma}[The variable layer bound]\label{lem:layers-variable}
The last two terms of Equation~\eqref{eq:21} satisfy
\begin{equation}
 -\mathcal L_I[H-P_0]-\mathcal L_T[H_0]
 \le\frac{\alpha(\log2+.25)}4\,\|[L,T]\|_2^2
       -\int(W-W_0)\,d\beta.
 \tag{22}\label{eq:22}
\end{equation}
\end{lemma}

\begin{proof}
Work in a spectral basis of \(L\), and use
\(a_i,\epsilon_i,E_z^L,S_z\) as in
Lemma~\ref{lem:layers-diagonal}.
Write \(g_i=g_z(x_i)\).
The kernel multiplying \(T_{ji}(E_z^L)_{ij}\) in
\(\mathcal L_T[H_0]\) is
\[
 \frac{(x_i-z)g_i+(x_j-z)g_j}{2}.
\]
Separate this kernel into the comparison term
\(\epsilon_i\sqrt{a_i a_j}(g_i+g_j)/2\) on a common-sign block,
zero on the other blocks, and a remainder \(k_z(x_i,x_j)\).
The original pairing therefore equals
\(\operatorname{tr}(S_z(T\circ g_z(L)))\) plus the pairing with
\(k_z\). We will bound the comparison term from below and the
absolute value of the remainder from above.

We first bound this pairing in matrix order.
For any real symmetric matrix \(V_0\), the spectral range of \(T\)
implies
\[
 T\circ V_0\ge
 t_cV_0-\frac{t_r}{2}\left(h_sI+\frac{V_0^2}{h_s}\right).
\]
Indeed, writing \(D=T-t_cI\), we have \(\|D\|\le t_r\), and for
every vector \(v_0\),
\[
 |\langle v_0,DV_0v_0\rangle|
 \le t_r|v_0|\,|V_0v_0|
 \le\frac{t_r}{2}
       \left(h_s|v_0|^2+h_s^{-1}|V_0v_0|^2\right).
\]
This proves the matrix inequality without a commutation
assumption or a sign restriction on \(V_0\).
Apply it with \(V_0=g_z(L)\), and pair against \(S_z\ge0\).
After integration, the comparison part of
\(\mathcal L_T[H_0]\) is bounded below by
\[
 \int\left[t_cg_z-\frac{t_r}{2}
                    (h_s+g_z^2/h_s)\right]\,d\beta.
\]
Together with
\(\mathcal L_I[H-P_0]=\int\overline{H-P_0}\,d\beta\),
this accounts for the term \(-\int W\,d\beta\).

It remains to estimate the pairing with the remainder \(k_z(x_i,x_j)\).
Use the nonnegative scalar weight
\[
 D_z=\frac{a_iW_0(z,x_i)+a_jW_0(z,x_j)}2.
\]
Lemma~\ref{lem:layers-diagonal} gives
\[
 \frac1m\mathbb E\int\sum_{i,j}
          D_z|(E_z^L)_{ij}|^2\,dz
 \le\int W_0\,d\beta.
\]
We will prove the scalar estimate
\[
 \int_{\mathbb R}\frac{k_z(x_i,x_j)^2}{D_z}\,dz
 \le\alpha(\log2+.25)(x_i-x_j)^2.
\]
Whenever \(D_z=0\), the defining integrals of \(H_0(z,\cdot)^2\)
on both relevant segments vanish.  The corresponding integrals
of \(H_0\), and hence \(k_z\), vanish as well; we set the quotient
equal to zero there.
For coincident endpoints, \(k_z\) vanishes identically.
We therefore assume below that \(\ell=|x_i-x_j|>0\).

If \(z\) lies between the endpoints, put
\(a=|x_i-z|\), \(b=|x_j-z|\), so \(a+b=\ell\).
Let \(Q_z\) be the integral of \(H_0(z,y)^2\) over the interval
between \(x_i,x_j\).
Then \(D_z=Q_z/(2\alpha)\).
Twice \(k_z\) is a signed sum of the integrals of \(H_0\)
over the two subintervals, so Cauchy--Schwarz gives
\[
 k_z^2\le\frac{\ell Q_z}{4},\qquad
 \frac{k_z^2}{D_z}\le\frac{\alpha\ell}{2}.
\]
The integral over this middle interval is therefore at most
\(\alpha\ell^2/2\).

If \(z\) lies outside the interval, order the distances so that
\(a\ge b>0\), and let \(\epsilon\) be their common sign.
Writing \(g_a,g_b\) for the two segment averages, we have
\[
 k_z^2
 =\frac{(\sqrt a-\sqrt b)^2}{4}
          (\sqrt a\,g_a-\sqrt b\,g_b)^2.
\]
Set \(h_z(t)=H_0(z,z+\epsilon t)\), for \(0\le t\le a\).
Then
\[
 \sqrt a\,g_a-\sqrt b\,g_b
 =\frac1{\sqrt a}\int_b^a h_z(t)\,dt
   +\left(\frac1{\sqrt a}-\frac1{\sqrt b}\right)
                         \int_0^b h_z(t)\,dt.
\]
Furthermore,
\[
 2\alpha D_z
 =\int_b^a h_z(t)^2\,dt+2\int_0^b h_z(t)^2\,dt.
\]
Cauchy--Schwarz, with multiplicity \(1\) on \((b,a]\)
and multiplicity \(2\) on \([0,b]\), consequently gives
\[
 (\sqrt a\,g_a-\sqrt b\,g_b)^2
 \le 2\alpha D_z
       \left[\frac{a-b}{a}
                    +\frac12(1-\sqrt{b/a})^2\right].
\]
In either exterior interval, write \(a=\ell(1+t)\),
\(b=\ell t\).  Integrating the last estimate over both
exterior intervals gives at most
\[
 \alpha\ell^2\int_0^\infty
  (\sqrt{1+t}-\sqrt t)^2
  \left[\frac1{1+t}
             +\frac12\left(1-\sqrt{\frac{t}{1+t}}\right)^2\right]dt.
\]
Use the substitution \(u_0=(\sqrt{1+t}-\sqrt t)^2\)
from the proof of Lemma~\ref{lem:layers-constant}.
The integral becomes
\[
 \int_0^1\left(\frac1{1+u_0}-\frac{u_0}{2}\right)du_0
 =\log2-\frac14.
\]
Adding the middle-interval contribution proves the scalar
estimate with coefficient
\(\alpha(\log2-\tfrac14+\tfrac12)
=\alpha(\log2+\tfrac14)\).

Put \(\mathcal V=\int W_0\,d\beta\).
Weighted Cauchy--Schwarz in \(G,z,i,j\), followed by the scalar
estimate, bounds the absolute value of the remaining pairing by
\[
 \sqrt{\mathcal V}\,
 \sqrt{\alpha(\log2+.25)}\,\|[L,T]\|_2.
\]
The convention for zero weights already covers \(\mathcal V=0\).
The comparison term has already contributed \(-\int W\,d\beta\)
to the upper bound for
\(-\mathcal L_I[H-P_0]-\mathcal L_T[H_0]\).
The remainder can increase this bound by at most the absolute value
just estimated. Writing \(W=(W-W_0)+W_0\), we obtain
\[
 -\int(W-W_0)\,d\beta
 -\mathcal V+\sqrt{\alpha(\log2+.25)\mathcal V}\,\|[L,T]\|_2.
\]
Completion of the square in \(\sqrt{\mathcal V}\) proves
Equation~\eqref{eq:22}.
\end{proof}

\subsection{The resulting bound on the linear field}

Applying Lemma~\ref{lem:layers-variable} in
Equation~\eqref{eq:21} leaves one term involving \(R\).
Since \(v\) is even and \(L\) is odd in \(G\), the field \(v'(L)\)
is odd.  The field \(R_o\) is orthogonal to degree one, so
\[
 \langle R,v'(L)\rangle_T
 =\langle R_o,v'(L)-L\rangle_T.
\]
The scalar estimate \(|v'(x)-x|\le.319\), applied by spectral
calculus, gives \(\|v'(L)-L\|\le.319\).
Hilbert--Schmidt Cauchy--Schwarz therefore yields
\[
 |\langle R_o,v'(L)-L\rangle_T|
 \le .319\,(\tau R_o^2)^{1/2}(\tau T^2)^{1/2}.
\]
For example, this follows directly from
\(\|T(v'(L)-L)\|_2^2\le .319^2\tau T^2\);
no commutation with \(T\) is needed.
As \(\Sigma\ge(1+t_-)I\), completing the square gives
\[
 -\eta\tau_\Sigma R_o^2-\langle R,v'(L)\rangle_T
 \le\frac{.319^2}{4\eta(1+t_-)}\,\tau T^2.
\]

For \(t\in[t_-,t_+]\), the scalar logarithm satisfies
\[
 t-\log(1+t)
 =t^2\int_0^1\frac{r_0}{1+r_0t}\,dr_0
 \le\frac{t^2}{2(1+t_-)}.
\]
Applying this to the eigenvalues of \(T\), we obtain
\[
 \begin{aligned}
 &-\eta\tau_\Sigma R_o^2-\langle R,v'(L)\rangle_T
                +\tau(T-\log\Sigma)\\
 &\hspace{8mm}\le
 \left(\frac{.319^2}{4\eta(1+t_-)}
                   +\frac1{2(1+t_-)}\right)\tau T^2
 =\frac{4417}{2608}\tau T^2
 \le1.8\,\tau T^2.
 \end{aligned}
\]
The last coefficient has margin
\(1.8-4417/2608=1387/13040>0\).
The commutator coefficient also has a positive margin:
\[
 e_0+\frac{\alpha(\log2+.25)}4
 <.256+\frac{2.26(.694+.25)}4
 =.78936<.80.
\]
Thus Equations~\eqref{eq:21} and~\eqref{eq:22} imply
\begin{equation}
 \mathcal E\le
 \Delta_{\rm p}-J_\Sigma(u)+.80\|[L,T]\|_2^2
       +1.8\,\tau T^2-\int(W-W_0)\,d\beta.
 \tag{23}\label{eq:23}
\end{equation}
Every term on the right except the final integral is now a
functional of the linear Gaussian field \(L\) and the deterministic
matrix \(T\).  The final integral is nonnegative before its minus
sign, because \(W>W_0\) and \(\beta\) is a positive measure.
